% !TeX program = pdflatex
\documentclass[11pt,a4paper]{article}

\usepackage[margin=2.5cm]{geometry}
\usepackage[T1]{fontenc}
\usepackage[utf8]{inputenc}
\usepackage{lmodern}
\usepackage{microtype}
\usepackage{booktabs}
\usepackage{tabularx}
\usepackage{enumitem}
\usepackage{xcolor}
\usepackage[colorlinks=true,linkcolor=blue,urlcolor=blue,citecolor=blue]{hyperref}

\newcolumntype{Y}{>{\raggedright\arraybackslash}X}

\title{\textbf{Evaluation Report}}
\author{gommage}
\date{June 2026}

\begin{document}
\maketitle

\begin{abstract}
This document reports the evaluation metrics for the gommage explainability
middleware. The evaluation compares agent execution before and after middleware
integration and measures response latency, token overhead for reasoning
provenance, and replay fidelity. The dataset was manually generated from 10
execution scenarios using \texttt{gpt-o4-mini} as the agent LLM with temperature
set to 0.85.
\end{abstract}

\section{Evaluation Setup}

The evaluation was designed to measure whether adding the AER-generating
middleware changes the practical behavior of the agent while still producing
useful reasoning provenance.

\begin{itemize}[leftmargin=1.5em]
  \item \textbf{Execution scenarios:} 10 manually generated agent execution scenarios.
  \item \textbf{LLM:} \texttt{gpt-o4-mini}.
  \item \textbf{Temperature:} 0.85.
  \item \textbf{Comparison:} before middleware integration versus after middleware integration.
  \item \textbf{Recorded artifact:} Agent Execution Record objects containing reasoning provenance.
\end{itemize}

\section{Metrics}

The evaluation focuses on three metrics:

\begin{enumerate}[leftmargin=1.5em]
  \item \textbf{Response latency:} the average change in response latency after
  middleware integration compared with the baseline agent execution.
  \item \textbf{Token overhead:} the average additional token cost required to
  generate AER reasoning provenance.
  \item \textbf{Replay fidelity:} the percentage of runs whose replay preserved
  the recorded execution trajectory across the evaluated scenarios.
\end{enumerate}

\section{Results}

\begin{table}[h]
\centering
\begin{tabularx}{\textwidth}{@{}p{0.28\textwidth}p{0.22\textwidth}Y@{}}
\toprule
\textbf{Metric} & \textbf{Result} & \textbf{Interpretation} \\
\midrule
Response latency delta &
6 ms average delta across 10 runs; $p = 0.24$ &
The measured latency difference before and after middleware integration was
not statistically significant. \\
Token overhead &
40 tokens on average &
Generating AER objects introduced a small average overhead for reasoning
provenance capture. \\
Replay fidelity &
100\% across 10 runs &
All evaluated executions replayed successfully against their recorded
trajectory. \\
\bottomrule
\end{tabularx}
\end{table}

\section{Latency Interpretation}

The average response latency delta was 6 ms across 10 runs. The associated
$p$-value was 0.24, which does not provide evidence of a statistically
significant latency difference between the baseline agent and the middleware
integrated agent. In this evaluation, the middleware did not introduce a
meaningful latency penalty.

\section{Token Overhead Interpretation}

The average token overhead for generating AER objects was 40 tokens. This
overhead corresponds to the additional structured reasoning provenance emitted
for each execution. The measured value indicates that the explainability layer
adds compact provenance metadata rather than a large auxiliary trace.

\section{Replay Fidelity Interpretation}

Replay fidelity was 100\% across the 10 manually generated scenarios. This means
that each recorded execution could be replayed from its AER without divergence
from the recorded trajectory under the evaluation procedure.

\section{Conclusion}

The evaluation shows that gommage provides reasoning provenance through AER
generation while preserving practical execution performance. The latency delta
was small and not statistically significant, the token overhead averaged 40
tokens, and replay fidelity reached 100\% across the evaluated scenarios.

\end{document}
